\documentclass[a4paper,11pt]{article}
\usepackage{exam-style-341}
\examinehead{341 农业综合知识三 · 数据库分章练习}{第04章\quad 规范化与模式设计}
\begin{document}

\unitcover{341 农业综合知识三 · 数据库分章练习}{第04章\quad 规范化与模式设计}{练习（共 100 分）}{约 60 分钟}{函数依赖（完全、部分、传递）；属性集闭包与候选码；第一范式、第二范式、第三范式和 BCNF；模式分解的无损连接性与保持函数依赖性；最小依赖集；关系模式分解方法；E-R 图向关系模式的转换与逻辑结构设计。}

\begin{knowmap}
\textbf{函数依赖：}$X\to Y$ 表示 $X$ 的取值决定 $Y$ 的取值；按依赖程度分为完全函数依赖、部分函数依赖和传递函数依赖\\[2pt]
\textbf{闭包与候选码：}$X^{+}$ 是由 $X$ 依据 $F$ 能推出的全部属性；若 $X^{+}$ 含全部属性且 $X$ 极小，则 $X$ 为候选码；求候选码先求单属性闭包，再逐步组合\\[2pt]
\textbf{范式阶梯：}1NF（属性原子）$\rightarrow$ 2NF（消除非主属性对码的部分依赖）$\rightarrow$ 3NF（消除非主属性对码的传递依赖）$\rightarrow$ BCNF（每个决定因素都包含候选码）\\[2pt]
\textbf{分解的两个标准：}无损连接性（分解后自然连接能还原原关系，不丢失也不多出元组）与保持函数依赖性（原 $F$ 中的依赖都能由分解后的模式蕴含）\\[2pt]
\textbf{无损连接判定：}二元分解 $\{R_1,R_2\}$ 无损的充要条件是 $R_1\cap R_2\to(R_1-R_2)$ 或 $R_1\cap R_2\to(R_2-R_1)$；多步分解可逐步判定或用 Chase 表格法\\[2pt]
\textbf{E-R 转换：}1:1 联系可合并到任一端；1:n 联系把 1 端主码并入 n 端；m:n 联系单独转换为关系模式，主码为两端主码的组合
\end{knowmap}

\begin{keyterms}
\textbf{函数依赖}：设 $R(U)$ 是属性集 $U$ 上的关系模式，$X,Y\subseteq U$，若对 $R$ 的任意两个元组，只要它们在 $X$ 上的值相同，在 $Y$ 上的值也必相同，则称 $X$ 函数决定 $Y$，记作 $X\to Y$。\\
\textbf{完全函数依赖}：若 $X\to Y$，且对 $X$ 的任何一个真子集 $X'$ 都有 $X'\not\to Y$，则称 $Y$ 完全函数依赖于 $X$。\\
\textbf{部分函数依赖}：若 $X\to Y$，但 $Y$ 不完全函数依赖于 $X$（即存在 $X$ 的真子集能决定 $Y$），则称 $Y$ 部分函数依赖于 $X$。\\
\textbf{传递函数依赖}：若 $X\to Y$、$Y\to Z$，且 $Y\not\to X$、$Z\not\subseteq Y$，则称 $Z$ 传递函数依赖于 $X$。\\
\textbf{属性集闭包}：属性集 $X$ 关于函数依赖集 $F$ 的闭包 $X^{+}$ 是由 $X$ 出发、依据 $F$ 能够推出的全部属性的集合。\\
\textbf{候选码}：能唯一标识元组且不含多余属性的属性组；用闭包判定，即 $K^{+}=U$ 且 $K$ 的任何真子集闭包都不等于 $U$。\\
\textbf{第一范式}：关系模式中每个属性都是不可再分的原子值，即满足 1NF。\\
\textbf{第二范式}：在 1NF 基础上，每个非主属性都完全函数依赖于候选码，即消除了非主属性对码的部分函数依赖。\\
\textbf{第三范式}：在 2NF 基础上，每个非主属性都不传递函数依赖于候选码，即消除了非主属性对码的传递函数依赖。\\
\textbf{BCNF}：在 3NF 基础上，每个决定因素都包含候选码（即所有非平凡函数依赖的左部都是超码）。\\
\textbf{无损连接分解}：关系模式 $R$ 分解为 $\rho=\{R_1,\dots,R_n\}$ 后，若对 $R$ 的任一合法关系 $r$ 都有 $r=\pi_{R_1}(r)\bowtie\cdots\bowtie\pi_{R_n}(r)$，则称该分解具有无损连接性。\\
\textbf{保持函数依赖}：若分解后各关系模式的函数依赖集的并集能够蕴含原函数依赖集 $F$ 中的每一个依赖，则称分解保持函数依赖。\\
\textbf{最小依赖集}：函数依赖集 $F$ 的最小依赖集 $F_{\min}$ 满足：每个依赖右部为单属性、左部没有多余属性、没有冗余的函数依赖。\\
\textbf{模式分解}：把一个关系模式分解为若干个关系模式，用以消除数据冗余和插入、删除、更新异常，同时尽量保证无损连接和保持函数依赖。
\end{keyterms}

\begin{contrast}{1NF / 2NF / 3NF / BCNF（每行依次为：必须满足的条件·消除的主要问题·判别要点）}
\citem{1NF}{每个属性都是不可再分的原子值}{属性可再分造成的表中有表}{检查每个属性是否原子}
\citem{2NF}{1NF 且非主属性完全函数依赖于候选码}{非主属性对码的部分函数依赖}{看候选码是否为组合码、非主属性是否只依赖码的一部分}
\citem{3NF}{2NF 且非主属性不传递函数依赖于候选码}{非主属性对码的传递函数依赖}{看是否存在“码 $\to$ 非主属性 $\to$ 非主属性”}
\citem{BCNF}{每个决定因素都包含候选码}{主属性对码的部分和传递函数依赖}{看每个非平凡函数依赖的左部是否为超码}
\end{contrast}

\begin{contrast}{完全函数依赖 / 部分函数依赖 / 传递函数依赖（每行依次为：定义·判定关键·典型例子）}
\citem{完全函数依赖}{$X\to Y$ 且 $X$ 的任何真子集都不能决定 $Y$}{依赖的是整个 $X$}{（学号，课程号）$\to$ 成绩}
\citem{部分函数依赖}{$X\to Y$，但 $X$ 的某个真子集也能决定 $Y$}{依赖 $X$ 的一部分}{（学号，课程号）$\to$ 姓名，而学号 $\to$ 姓名}
\citem{传递函数依赖}{$X\to Y$、$Y\to Z$，且 $Y\not\to X$、$Z\not\subseteq Y$}{经过中间属性间接决定}{学号 $\to$ 系号，系号 $\to$ 系名}
\end{contrast}

\begin{contrast}{无损连接分解 / 保持函数依赖分解（每行依次为：要解决的问题·判定方法·二者的关系）}
\citem{无损连接分解}{分解后自然连接不丢失、不多出元组}{二元分解定理；Chase 表格法}{与保持函数依赖相互独立，可同时满足，也可只满足其一}
\citem{保持函数依赖分解}{分解后原函数依赖不被破坏}{检查 $F$ 中每个依赖能否由分解后的模式蕴含}{满足保持函数依赖的 3NF 分解一定存在，但不一定达到 BCNF}
\end{contrast}

\qtypename{一、单项选择题}{10 题，每题 2 分，共 20 分}
\q 关系模式满足第一范式（1NF）的条件是（\quad）
\choices{每个属性都是不可再分的原子值}{每个非主属性都完全函数依赖于码}{每个非主属性都不传递函数依赖于码}{每个决定因素都包含候选码}

\q 第二范式（2NF）是在 1NF 的基础上消除了（\quad）
\choices{非主属性对码的部分函数依赖}{非主属性对码的传递函数依赖}{主属性对码的部分函数依赖}{所有函数依赖}

\q 第三范式（3NF）要求关系模式（\quad）
\choices{每个非主属性都完全函数依赖于码}{每个非主属性都不传递函数依赖于码}{每个决定因素都包含候选码}{属性不可再分}

\q BCNF 要求关系模式中的每个非平凡函数依赖（\quad）
\choices{左部都是单属性}{右部都是单属性}{左部都包含候选码}{右部都包含候选码}

\q 任何二元关系模式最高可以达到的范式是（\quad）
\choicesline{A. 1NF}{B. 2NF}{C. 3NF}{D. BCNF}

\q 设关系模式 $R$ 分解为 $\{R_1,R_2\}$，则判定该分解具有无损连接性的常用条件是（\quad）
\choices{$R_1\cap R_2\to(R_1-R_2)$ 或 $R_1\cap R_2\to(R_2-R_1)$}{$R_1\cup R_2=R$}{$R_1\cap R_2=\varnothing$}{$R_1-R_2\to R_2$}

\q 设关系模式 $R(A,B,C,D)$，$F=\{A\to B,\ B\to C,\ C\to D\}$，则 $A^{+}$ 为（\quad）
\choicesline{A. $AB$}{B. $ABC$}{C. $ABCD$}{D. $AD$}

\q 函数依赖集 $F$ 的最小依赖集 $F_{\min}$ 不必满足的条件是（\quad）
\choices{每个函数依赖的右部为单属性}{每个函数依赖的左部没有多余属性}{没有冗余的函数依赖}{每个函数依赖的左部都是单属性}

\q 若要求模式分解既保持函数依赖又具有无损连接性，则分解至少可以达到（\quad）
\choicesline{A. 1NF}{B. 2NF}{C. 3NF}{D. BCNF}

\q 下列关于模式分解的说法中，正确的是（\quad）
\choices{保持函数依赖的分解一定具有无损连接性}{无损连接的分解一定保持函数依赖}{保持函数依赖的分解一定能达到 BCNF}{无损连接性和保持函数依赖是两个相互独立的分解标准}

\qtypename{二、多项选择题}{5 题，每题 4 分，共 20 分；每题有两个或两个以上正确项}
\q 按照函数依赖的程度，函数依赖可以分为（\quad）
\choicesline{A. 完全函数依赖}{B. 部分函数依赖}{C. 传递函数依赖}{D. 平凡函数依赖}

\q 关于范式的下列说法中，正确的有（\quad）
\choices{1NF 要求属性不可再分}{2NF 消除了非主属性对码的部分函数依赖}{3NF 消除了非主属性对码的传递函数依赖}{BCNF 消除了所有函数依赖}

\q 判定模式分解具有无损连接性的方法有（\quad）
\choices{二元分解判定定理}{Chase 表格法}{检查属性集闭包}{检查分解是否保持函数依赖}

\q 函数依赖集 $F$ 的最小依赖集应满足（\quad）
\choices{每个函数依赖的右部都是单属性}{每个函数依赖的左部没有多余属性}{没有冗余的函数依赖}{所有函数依赖的左部都必须是单属性}

\q 将 E-R 图转换为关系模式时，下列说法正确的有（\quad）
\choices{每个实体转换为一个关系模式}{1:n 联系把 1 端实体的主码作为外码加入 n 端}{m:n 联系单独转换为一个关系模式}{1:1 联系只能单独转换为一个关系模式}

\qtypename{三、判断题}{5 题，每题 2 分，共 10 分；正确的画“√”，错误的画“×”并写出改正后的表述}
\q 任何二元关系模式都属于 BCNF。\quad（\quad）\quad 改正：\fixblank

\q 满足 3NF 的关系模式一定满足 BCNF。\quad（\quad）\quad 改正：\fixblank

\q 保持函数依赖的分解一定具有无损连接性。\quad（\quad）\quad 改正：\fixblank

\q 若 $X\to Y$ 且 $Y\to Z$，则 $X\to Z$。\quad（\quad）\quad 改正：\fixblank

\q BCNF 在 3NF 的基础上进一步消除了主属性对码的部分函数依赖和传递函数依赖。\quad（\quad）\quad 改正：\fixblank

\qtypename{四、填空题}{5 题，每题 2 分，共 10 分}
\q 关系模式的范式级别由低到高依次为 1NF、\fillblank{}、\fillblank{} 和 BCNF。

\q 2NF 消除了非主属性对码的 \fillblank{}；3NF 消除了非主属性对码的 \fillblank{}。

\q 设关系模式 $R(A,B,C,D)$，$F=\{A\to B,\ B\to C,\ C\to D\}$，则 $A^{+}=$ \fillblank{}，$R$ 的候选码是 \fillblank{}。

\q 二元分解 $\{R_1,R_2\}$ 具有无损连接性的充要条件是 $R_1\cap R_2\to(R_1-R_2)$ 或 \fillblank{}。

\q 最小依赖集要求每个函数依赖的右部为单属性、左部没有多余属性，并且 \fillblank{}。

\qtypename{五、简答题}{2 题，每题 8 分，共 16 分}
\q 试述 1NF、2NF、3NF 和 BCNF 的定义，并说明它们之间的包含关系。\anslines{2.2cm}

\q 什么是无损连接分解？什么是保持函数依赖分解？为什么关系模式分解要同时考虑这两个标准？\anslines{2.2cm}

\qtypename{六、综合题}{2 题，每题 12 分，共 24 分}
\q 设关系模式 $R(\text{学号},\text{课程号},\text{教师号},\text{教师名},\text{教室})$，简记为 $R(\text{Sno},\text{Cno},\text{Tno},\text{Tname},\text{Room})$，其函数依赖集为
\[
F=\{(\text{Sno},\text{Cno})\to\text{Tno},\quad \text{Tno}\to\text{Tname},\quad \text{Tno}\to\text{Room}\}.
\]
请完成：
\begin{enumerate}[label=（\arabic*）, leftmargin=2.4em, itemsep=1pt]
  \item 求 $R$ 的全部候选码，并指出主属性和非主属性。（3 分）
  \item $R$ 最高属于第几范式？说明理由。（3 分）
  \item 求 $F$ 的最小依赖集 $F_{\min}$。（2 分）
  \item 将 $R$ 分解为 3NF，要求分解既保持函数依赖又具有无损连接性，并逐步验证无损连接性。（4 分）
\end{enumerate}
\anslines{4.5cm}

\q 某高校教务管理系统的 E-R 图包含如下信息：
\begin{itemize}[leftmargin=1.6em, itemsep=1pt]
  \item 实体“系”：系号、系名；
  \item 实体“教师”：教师号、姓名、职称；
  \item 实体“课程”：课程号、课程名、学分；
  \item 实体“学生”：学号、姓名、性别；
  \item 一个系有多名教师，一名教师只属于一个系；
  \item 一名教师可以讲授多门课程，一门课程也可以由多名教师讲授，授课联系上记录学时；
  \item 一名学生可以选修多门课程，一门课程也可以被多名学生选修，选修联系上记录成绩。
\end{itemize}

请完成：
\begin{enumerate}[label=（\arabic*）, leftmargin=2.4em, itemsep=1pt]
  \item 指出上述三个联系的类型（1:1、1:n 或 m:n）。（3 分）
  \item 说明 E-R 图转换为关系模式的基本规则。（3 分）
  \item 写出转换后得到的全部关系模式，标出每个关系模式的主码和外码。（6 分）
\end{enumerate}
\anslines{4.5cm}

\end{document}
